% spectral_residual_access.tex - isolated mathematical paper; depends on ../report-source.md and ../notes/minimax_residual_access.md.
\documentclass[11pt]{article}

\usepackage[a4paper,margin=28mm]{geometry}
\usepackage{amsmath,amssymb,amsthm,mathtools}
\usepackage{booktabs}
\usepackage{microtype}
\usepackage{enumitem}
\usepackage{xcolor}
\usepackage[hidelinks]{hyperref}

\definecolor{ink}{HTML}{17212B}
\definecolor{muted}{HTML}{566573}
\hypersetup{
  colorlinks=true,
  linkcolor=ink,
  citecolor=ink,
  urlcolor=muted,
  pdftitle={Spectral Residual-Access Complexity},
  pdfauthor={Cassette Pure Mathematics Laboratory},
  pdfsubject={Exact minimax laws for query-independent categorical matrix execution}
}

\newtheorem{theorem}{Theorem}[section]
\newtheorem{proposition}[theorem]{Proposition}
\newtheorem{corollary}[theorem]{Corollary}
\newtheorem{lemma}[theorem]{Lemma}
\theoremstyle{definition}
\newtheorem{definition}[theorem]{Definition}
\newtheorem{remark}[theorem]{Remark}

\newcommand{\F}{\mathbb F}
\newcommand{\D}{\mathcal D}
\newcommand{\tr}{\operatorname{tr}}
\newcommand{\diag}{\operatorname{diag}}
\newcommand{\ran}{\operatorname{ran}}
\newcommand{\rank}{\operatorname{rank}}
\newcommand{\E}{\mathbb E}
\newcommand{\Var}{\operatorname{Var}}
\newcommand{\norm}[1]{\left\lVert #1\right\rVert}
\newcommand{\inner}[2]{\left\langle #1,#2\right\rangle}

\title{\textbf{Spectral Residual-Access Complexity}\\
\large Exact Minimax Laws for Query-Independent Categorical Matrix Execution}
\author{Cassette Pure Mathematics Laboratory\\
\small Isolated research manuscript}
\date{5 September 2026}

\begin{document}
\maketitle

\begin{abstract}
Let a residual operator be evaluated by drawing one of its columns, or one of its physically
fixed page components, from a probability law chosen before the query is known.  The present
squared-energy law supplies a simple Frobenius-scale certificate.  It need not minimize the
largest query variance.  We determine the latter problem exactly for independent categorical
sampling with inverse-probability weights.

For a matrix $R$ with nonzero columns $r_j$, column energies $a_j=\norm{r_j}_2^2$, and
$G=R^*R$, the one-access worst-query value is
\[
  \nu(R)=\min_{\pi\in\Delta_q^\circ}
  \lambda_{\max}\!\left(\diag(a_j/\pi_j)-G\right).
\]
We prove the density-matrix dual
\[
  \nu(R)=\max_{X\succeq0,\,\tr X=1}
  \left\{\left(\sum_j\sqrt{a_jX_{jj}}\right)^2-\tr(XG)\right\},
\]
derive a finite optimality certificate that remains valid when the largest eigenvalue is
multiple, and obtain exact laws for two columns, orthogonal columns, and every rank-one residual.
We also characterize precisely when squared-energy sampling is minimax.  For two orthogonal
columns, its exact variance can exceed the minimax value by an arbitrarily large factor even at
fixed Frobenius norm.  The same duality extends to a fixed decomposition into physical page
operators.

Inverse-probability estimation and E-optimal covariance design are prior art.  The contribution
claimed here is the exact combined analysis of this reciprocal, fixed-law access problem.  A
bounded three-wave primary-source search found no earlier source stating that combination; such a
search cannot establish literature-wide novelty.  This manuscript is pure mathematical research.
It does not amend Cassette's production mathematics or implementation.
\end{abstract}

\section{Introduction}

Suppose that a linear residual $R$ is too expensive to keep wholly resident.  One elementary
unbiased procedure stores an access law, draws a component, and rescales the component by the
inverse of its probability.  If the query is already known, the probability can follow that
query.  Our concern is harder and more rigid: one law must be fixed before any query arrives.
The relevant loss is then the largest mean-square error over the unit sphere.

This distinction changes the geometry.  Squared-column-norm sampling makes a simple Frobenius
bound available, but that sufficient rule need not be spectrally optimal.  The exact object is not
$\norm{R}_F^2$ alone.  It is a coordinate-sensitive minimax invariant, denoted $\nu(R)$ below,
which remembers both column energies and their Gram coherence.

The main results are as follows.
\begin{enumerate}[label=\textup{(\roman*)},leftmargin=2.1em]
  \item The exact covariance is a reciprocal diagonal perturbation of $R^*R$, and $s$
  independent accesses divide the one-access minimax value by exactly $s$.
  \item The primal problem has an exact dual over density matrices.  Its saddle point gives a
  global certificate, including at repeated top eigenvalues where a single leading eigenvector
  is insufficient.
  \item Two-column residuals have a square-root law; orthogonal residuals have a reciprocal
  water-filling law.
  \item The squared-energy law is minimax precisely when its energy vector is the diagonal of a
  density matrix supported on the minimum eigenspace of $R^*R$.
  \item Rank-one residuals undergo a complete balanced--dominant phase transition governed by
  the polygon inequality.
  \item A normalized two-column family gives an unbounded multiplicative separation from the
  squared-energy law.
  \item The covariance-duality framework extends from columns to a physically fixed page
  decomposition, with a sharp zero-variance characterization.
\end{enumerate}

Theorems here are exact only inside the declared information class: fresh independent
categorical draws, ordinary averaging, inverse-probability weights, and a law independent of the
query.  Adaptive access, sampling without replacement, query-dependent laws, and arbitrary
randomized linear information are different problems.

\section{Prior work and the claim boundary}

Unequal-probability sampling with replacement and inverse-probability weighting go back at least
to Hansen and Hurwitz~\cite{HansenHurwitz1943}.  Categorical matrix-product estimators, including
norm-product probabilities that minimize expected Frobenius error, appear in the randomized
linear-algebra literature~\cite[\S4.1, Lemmas 3--4]{DrineasKannanMahoney2006}.  Query matching with probabilities
chosen for an expected-variance objective is studied by Eriksson-Bique et
al.~\cite[\S\S2--4]{ErikssonBiqueEtAl2011}.  Spectral random approximation of Gram matrices and entrywise
sparsification address related but different random objects and losses
\cite[\S\S3--5]{HolodnakIpsen2015}; see also the contextual comparison with entrywise
sparsification in~\cite{AchlioptasMcSherry2007}.

The closest broad precedent found in the recorded search is the E-optimal inverse-probability
subsampling formulation of Imberg, Axelson-Fisk, and Jonasson
\cite[eqs.~(8)--(9), Table~1]{ImbergEtAl2023}.  Their multinomial covariance and
largest-eigenvalue objective include the ancestry of our problem; their simple-eigenvalue
stationarity law also anticipates the square-root form of a differentiable optimum
\cite[Lemma~2(c), Proposition~1, eq.~(13)]{ImbergEtAl2023}.  Classical general equivalence theory
frames E-optimal design~\cite{Kiefer1974}; a modern treatment explicitly uses trace-one positive
semidefinite witnesses on an extremal eigenspace~\cite[\S2]{HarmanRosa2018}.  Sion's theorem supplies the minimax interchange used
here~\cite{Sion1958}.  Schur--Horn geometry and prescribed-norm tight frames supply nearby
diagonal and balancing ideas~\cite{Horn1954,CasazzaEtAl2006}.  In particular, the latter records
the prescribed-norm fundamental inequality~\cite[Proposition~1]{CasazzaEtAl2006}.  Novak's
randomized Hilbert-space operator problem uses a different information class and exhibits related allocation geometry
\cite{Novak1992}.  Specialized exact E-optimal allocations also occur in factorial design
\cite[Theorem~1]{RavichandranEtAl2024}.

Our claim is deliberately narrower than these subjects.  Within the three-wave primary-source
search recorded with this manuscript, with cutoff 5 September 2026, we found no earlier source
combining a fixed, query-independent, categorical with-replacement residual-access law with the
reciprocal density-matrix dual, a repeated-top-eigenvalue certificate, the closed forms proved
below, and the exact squared-energy optimality criterion.  This is evidence from a finite search,
not proof of universal novelty.  In particular, we do not claim to invent inverse-probability
sampling, E-optimal covariance design, or density witnesses for extremal eigenvalues.

\section{The categorical access model}

Let $\F$ be $\mathbb R$ or $\mathbb C$.  Let
$R:\F^q\to\F^p$ be nonzero, with columns $r_1,\ldots,r_q$.  Zero columns are discarded: they
require neither sampling mass nor residual access.  Put
\[
  G=R^*R,\qquad a_j=G_{jj}=\norm{r_j}_2^2>0,
\]
and let
\[
  \Delta_q^\circ=\left\{\pi\in\mathbb R^q:
  \pi_j>0,\ \sum_{j=1}^q\pi_j=1\right\}.
\]
For $\pi\in\Delta_q^\circ$, draw $J$ with $\Pr(J=j)=\pi_j$.  Given a query
$x\in\F^q$, return
\begin{equation}
  Z_\pi(x)=\frac{r_Jx_J}{\pi_J}.
  \label{eq:estimator}
\end{equation}
The law $\pi$ is fixed independently of $x$.

\begin{proposition}[Exact covariance]
\label{prop:covariance}
The estimator in \eqref{eq:estimator} is unbiased and
\begin{equation}
  \E\norm{Z_\pi(x)-Rx}_2^2=x^*Q_\pi x,
  \qquad
  Q_\pi=\diag\!\left(\frac{a_j}{\pi_j}\right)-G\succeq0.
  \label{eq:covariance}
\end{equation}
\end{proposition}

\begin{proof}
Since $Rx=\sum_jr_jx_j$,
\[
  \E Z_\pi(x)=\sum_j\pi_j\frac{r_jx_j}{\pi_j}=Rx.
\]
Also
\[
  \E\norm{Z_\pi(x)}_2^2
  =\sum_j\pi_j\frac{\norm{r_j}_2^2|x_j|^2}{\pi_j^2}
  =x^*\diag(a_j/\pi_j)x.
\]
Subtracting $\norm{Rx}_2^2=x^*Gx$ proves the identity.  Its left-hand side is nonnegative for
every $x$, hence $Q_\pi\succeq0$.
\end{proof}

\begin{definition}[Spectral residual-access complexity]
\label{def:nu}
The one-access fixed-law minimax value is
\begin{equation}
  \nu(R)=\min_{\pi\in\Delta_q^\circ}\lambda_{\max}(Q_\pi).
  \label{eq:primal}
\end{equation}
\end{definition}

The minimum in \eqref{eq:primal} is attained.  Indeed, the objective is continuous in the open
simplex, while for every $j$,
\[
  \lambda_{\max}(Q_\pi)\ge e_j^*Q_\pi e_j
  =a_j(\pi_j^{-1}-1)\longrightarrow\infty
  \quad\text{as }\pi_j\downarrow0.
\]
Thus every minimizing sequence stays in a compact subset of $\Delta_q^\circ$.

\begin{theorem}[Exact minimax access law]
\label{thm:sample-complexity}
Let $\overline Z_{\pi,s}$ be the arithmetic mean of $s\ge1$ independent copies of
$Z_\pi$.  Then
\begin{equation}
  \min_{\pi\in\Delta_q^\circ}\ 
  \sup_{\norm{x}_2=1}
  \E\norm{\overline Z_{\pi,s}(x)-Rx}_2^2
  =\frac{\nu(R)}{s}.
  \label{eq:s-samples}
\end{equation}
If $R\ne0$, $\varepsilon>0$, and $A\ne0$, the least integer $s\ge1$ for which the left side is
at most $\varepsilon^2\norm{A}_F^2$ is
\begin{equation}
  s_{\min}=\max\left\{1,
  \left\lceil\frac{\nu(R)}{\varepsilon^2\norm{A}_F^2}\right\rceil\right\}.
  \label{eq:sample-count}
\end{equation}
If $R=0$, zero residual accesses suffice.
\end{theorem}

\begin{proof}
The errors of the $s$ copies are independent and centered, so all cross terms vanish and the
covariance form in Proposition~\ref{prop:covariance} is divided by $s$.  Rayleigh--Ritz gives
$\sup_{\norm{x}=1}x^*Q_\pi x=\lambda_{\max}(Q_\pi)$.  Minimizing over $\pi$ proves
\eqref{eq:s-samples}, and the integer inequality gives \eqref{eq:sample-count}.  A nonzero
residual with $\nu(R)=0$ still needs one exact component access; zero variance does not make the
residual resident.
\end{proof}

\section{An exact density-matrix dual}

Let
\[
  \D_q=\{X\in\F^{q\times q}:X\succeq0,\ \tr X=1\}.
\]
This compact convex set replaces a single leading eigenvector by every convex mixture supported
on a leading eigenspace.

\begin{theorem}[Reciprocal density dual]
\label{thm:dual}
For every nonzero $R$ with zero columns removed,
\begin{equation}
  \boxed{
  \nu(R)=\max_{X\in\D_q}
  \left[
    \left(\sum_{j=1}^q\sqrt{a_jX_{jj}}\right)^2-\tr(XG)
  \right].}
  \label{eq:dual}
\end{equation}
\end{theorem}

\begin{proof}
Rayleigh--Ritz in density-matrix form states that, for Hermitian $H$,
\begin{equation}
  \lambda_{\max}(H)=\max_{X\in\D_q}\tr(XH).
  \label{eq:density-rayleigh}
\end{equation}
Define
\[
  L(\pi,X)=\sum_j\frac{a_jX_{jj}}{\pi_j}-\tr(XG).
\]
For $0<\rho<1/q$, let
$\Delta_{q,\rho}=\{\pi:\pi_j\ge\rho,\ \sum_j\pi_j=1\}$.  On
$\Delta_{q,\rho}\times\D_q$, the function $L$ is continuous, convex in $\pi$, and affine in
$X$.  Sion's minimax theorem therefore gives
\begin{equation}
  \min_{\pi\in\Delta_{q,\rho}}\max_{X\in\D_q}L(\pi,X)
  =\max_{X\in\D_q}\min_{\pi\in\Delta_{q,\rho}}L(\pi,X).
  \label{eq:sion}
\end{equation}
The coercivity noted after Definition~\ref{def:nu} places a primal minimizer in the interior;
hence the left side of \eqref{eq:sion} equals $\nu(R)$ for all sufficiently small $\rho$.

Fix $X\in\D_q$ and write $c_j=a_jX_{jj}\ge0$.  Cauchy--Schwarz gives
\[
  \left(\sum_j\sqrt{c_j}\right)^2
  =\left(\sum_j\sqrt{c_j/\pi_j}\sqrt{\pi_j}\right)^2
  \le\sum_j\frac{c_j}{\pi_j}.
\]
Equality is approached, and is attained when every $c_j>0$, by
$\pi_j=\sqrt{c_j}/\sum_k\sqrt{c_k}$.  Consequently
\begin{equation}
  \inf_{\pi\in\Delta_q^\circ}\sum_j\frac{c_j}{\pi_j}
  =\left(\sum_j\sqrt{c_j}\right)^2.
  \label{eq:cs-inner}
\end{equation}
To pass from the restricted simplexes to the open simplex without a hidden boundary step, choose
any sequence $\rho_n\downarrow0$.  The restricted inner-value functions are continuous on the
compact set $\D_q$, decrease pointwise to the continuous right side of
\eqref{eq:cs-inner} minus $\tr(XG)$, and therefore converge uniformly by Dini's theorem.  Their
maxima converge to the maximum in \eqref{eq:dual}.  Equation~\eqref{eq:sion} now proves the
claim.
\end{proof}

\begin{theorem}[Repeated-eigenvalue-safe certificate]
\label{thm:certificate}
A probability vector $\pi\in\Delta_q^\circ$ is minimax with value $t$ if and only if there is
$X\in\D_q$ with positive diagonal such that
\begin{align}
  Q_\pi&\preceq tI,                                      \label{eq:cert-upper}\\
  X(tI-Q_\pi)&=0,                                        \label{eq:cert-support}\\
  \pi_j&=\frac{\sqrt{a_jX_{jj}}}
  {\sum_k\sqrt{a_kX_{kk}}}\quad(1\le j\le q).           \label{eq:cert-law}
\end{align}
\end{theorem}

\begin{proof}
Assume first that the displayed conditions hold.  The first gives the primal bound
$\nu(R)\le t$.  Taking the trace in \eqref{eq:cert-support} gives
$\tr(XQ_\pi)=t$.  By \eqref{eq:cert-law},
\[
  \sum_j\frac{a_jX_{jj}}{\pi_j}
  =\left(\sum_j\sqrt{a_jX_{jj}}\right)^2.
\]
Thus the dual objective at $X$ equals $t$.  Weak duality gives $\nu(R)\ge t$, so equality holds
and $\pi$ is minimax.

Conversely, take a primal minimizer.  For sufficiently small $\rho$, it is also optimal in the
compact game \eqref{eq:sion}, which has a saddle point $(\pi,X)$.  Maximization in $X$ places
$\ran X$ in the top eigenspace of $Q_\pi$, yielding \eqref{eq:cert-upper} and
\eqref{eq:cert-support}.  Interior stationarity in $\pi$ gives
$a_jX_{jj}/\pi_j^2$ independent of $j$.  This common value cannot vanish because
$\tr X=1$ and every $a_j>0$; hence every $X_{jj}>0$, and normalization gives
\eqref{eq:cert-law}.
\end{proof}

The certificate is finite and exact.  A rank-one matrix $X=vv^*$ suffices when a simple top
eigenvalue exposes the optimum.  When the top eigenvalue is repeated, the required $X$ may be a
nontrivial mixture; using a single arbitrarily selected eigenvector can then miss the saddle
condition.

\begin{proposition}[Exact conic representation]
\label{prop:conic}
The dual is an exact semidefinite/second-order-cone optimization problem.  Let
$A_j=a_je_je_j^*$ and $B=G$.  Introduce
$b_j=\tr(XA_j)$ and nonnegative $z_{jk}$ subject to
$z_{jk}^2\le b_jb_k$.  Then \eqref{eq:dual} equals the maximum of
\begin{equation}
  \sum_jb_j+2\sum_{j<k}z_{jk}-\tr(XB)
  \label{eq:conic-objective}
\end{equation}
over $X\succeq0$, $\tr X=1$, the defining equations for $b_j$, and the stated rotated-cone
constraints.
\end{proposition}

\begin{proof}
For fixed nonnegative $b_j$, every coefficient of $z_{jk}$ in
\eqref{eq:conic-objective} is positive.  Therefore every maximizing solution has
$z_{jk}=\sqrt{b_jb_k}$.  The objective becomes
$(\sum_j\sqrt{b_j})^2-\tr(XB)$, exactly \eqref{eq:dual}.  No relaxation has been introduced.
\end{proof}

\section{Closed forms}

\subsection{Two columns}

\begin{theorem}[Two-column law]
\label{thm:two-column}
Let $R$ have two nonzero columns, and put
\[
  a=\norm{r_1}_2^2,\qquad b=\norm{r_2}_2^2,
  \qquad c=\inner{r_1}{r_2}.
\]
Then the unique minimax law and value are
\begin{equation}
  \pi^*=\frac{(\sqrt a,\sqrt b)}{\sqrt a+\sqrt b},
  \qquad
  \nu(R)=\sqrt{ab}+|c|.
  \label{eq:two-column}
\end{equation}
\end{theorem}

\begin{proof}
Write $u=\pi_1/(1-\pi_1)>0$.  Directly from \eqref{eq:covariance},
\[
  Q_\pi=
  \begin{pmatrix}
    a/u&-c\\
    -\overline c&bu
  \end{pmatrix}.
\]
Its determinant is $ab-|c|^2$, independent of $u$, while its trace is $a/u+bu$.  For a
positive semidefinite $2\times2$ matrix with fixed nonnegative determinant, the larger eigenvalue
is strictly increasing in the trace on the feasible interval, and here its minimum is unique because
\[
  a/u+bu\ge2\sqrt{ab},
\]
with equality only at $u=\sqrt{a/b}$.  At that point the eigenvalues are
$\sqrt{ab}\pm|c|$.  The larger one and the corresponding probability give
\eqref{eq:two-column}.
\end{proof}

The probability depends only on the two column norms.  Their coherence enters the value but not
the law.  This separation is special to two columns.

\subsection{Orthogonal columns}

\begin{theorem}[Reciprocal water filling]
\label{thm:orthogonal}
If the nonzero columns of $R$ are pairwise orthogonal, there is a unique $u\ge0$ satisfying
\begin{equation}
  \sum_{j=1}^q\frac{a_j}{a_j+u}=1.
  \label{eq:water-equation}
\end{equation}
The minimax value is $\nu(R)=u$, and the unique minimax law is
\begin{equation}
  \pi_j^*=\frac{a_j}{a_j+u}.
  \label{eq:water-law}
\end{equation}
For $q=1$, these formulas mean $u=0$ and $\pi_1=1$.
\end{theorem}

\begin{proof}
Now $G=\diag(a_j)$ and
\[
  Q_\pi=\diag\bigl(a_j(\pi_j^{-1}-1)\bigr).
\]
The inequality $Q_\pi\preceq uI$ is equivalent to
$\pi_j\ge a_j/(a_j+u)$ for every $j$.  A probability vector satisfying these lower bounds
exists precisely when their sum is at most one.  If $q\ge2$, that sum is continuous and strictly
decreasing from $q$ to $0$ as $u$ runs from $0$ to infinity, so the least feasible $u$ is the
unique solution of \eqref{eq:water-equation}.  At the least value, the lower bounds sum to one
and force \eqref{eq:water-law}; hence the law is unique.  The one-column boundary is immediate.
\end{proof}

\section{When squared-energy sampling is exact}

Let
\[
  T=\tr G=\sum_ja_j=\norm{R}_F^2,
  \qquad \pi_j^{\mathrm F}=\frac{a_j}{T}.
\]
This is the squared-column-norm, or energy, law.  It makes
\begin{equation}
  Q_{\pi^{\mathrm F}}=TI-G.
  \label{eq:frob-covariance}
\end{equation}
Let $E_{\min}(G)$ denote the eigenspace belonging to $\lambda_{\min}(G)$.

\begin{theorem}[Exact energy-law criterion]
\label{thm:energy-criterion}
The squared-energy law $\pi^{\mathrm F}$ is minimax if and only if there exists $X\in\D_q$ such
that
\begin{equation}
  \ran X\subseteq E_{\min}(G),
  \qquad
  \diag X=\left(\frac{a_1}{T},\ldots,\frac{a_q}{T}\right).
  \label{eq:energy-criterion}
\end{equation}
When this condition holds,
\begin{equation}
  \nu(R)=T-\lambda_{\min}(G).
  \label{eq:energy-value}
\end{equation}
When it fails, the minimax law has strictly smaller worst-query variance than the energy law.
\end{theorem}

\begin{proof}
Equation~\eqref{eq:frob-covariance} shows that its largest eigenvalue is the right side of
\eqref{eq:energy-value} and its top eigenspace is $E_{\min}(G)$.  Suppose first that $X$ satisfies
\eqref{eq:energy-criterion}.  Then
\[
  \frac{\sqrt{a_jX_{jj}}}{\sum_k\sqrt{a_kX_{kk}}}
  =\frac{a_j/\sqrt T}{\sqrt T}=\frac{a_j}{T}.
\]
Together with the support condition, Theorem~\ref{thm:certificate} proves minimaxity and
\eqref{eq:energy-value}.

Conversely, suppose $\pi^{\mathrm F}$ is minimax.  The necessary half of
Theorem~\ref{thm:certificate} supplies $X\in\D_q$ supported on the top eigenspace of
$Q_{\pi^{\mathrm F}}$, namely $E_{\min}(G)$.  Its stationarity equation says
\[
  \frac{a_jX_{jj}}{(a_j/T)^2}=\gamma
  \quad\text{for every }j
\]
for some common $\gamma>0$.  Hence $X_{jj}=\gamma a_j/T^2$.  Summing the diagonal and using
$\tr X=1$ yields $\gamma=T$, which is precisely \eqref{eq:energy-criterion}.  Thus failure of
that condition means the displayed energy law is not a minimizer; the minimum is then strictly
below its value.
\end{proof}

The criterion separates two questions that a Frobenius norm merges.  The vector
$(a_j/T)_j$ records marginal energy.  The minimum eigenspace records how the columns can cancel
under the worst query.  Energy sampling is exact only when one density matrix realizes both at
once.

\section{The complete rank-one phase transition}

Assume now that $\rank R=1$.  Then $G=zz^*$ for some $z\in\F^q$, with
$|z_j|^2=a_j$.  Put $T=\sum_ja_j$, let $a=\max_ja_j$, and put $b=T-a$.

\begin{theorem}[Rank-one law]
\label{thm:rank-one}
If $R$ has at least two nonzero columns, then
\begin{equation}
  \nu(R)=
  \begin{cases}
    T,&a\le T/2,\\[2mm]
    2\sqrt{ab},&a>T/2.
  \end{cases}
  \label{eq:rank-one-value}
\end{equation}
In the balanced regime $a\le T/2$, the unique minimax law is
$\pi_j^*=a_j/T$.  In the dominant regime $a>T/2$, let $h$ be the unique index with $a_h=a$.
Then the unique minimax law is
\begin{equation}
  \pi_h^*=\frac{\sqrt a}{\sqrt a+\sqrt b},
  \qquad
  \pi_j^*=\frac{a_j}{\sqrt b(\sqrt a+\sqrt b)}\quad(j\ne h).
  \label{eq:rank-one-law}
\end{equation}
For a single nonzero column, $\nu(R)=0$ and the sole probability is one.
\end{theorem}

\begin{proof}
First suppose $a\le T/2$.  We construct a certificate for the energy law.  The polygon
inequality says that planar vectors of prescribed lengths $a_j/\sqrt T$ can sum to zero exactly
when their largest length does not exceed half their total length.  Choose unit vectors $w_j$ in
a real plane so that
\[
  \sum_j\frac{a_j}{\sqrt T}w_j=0.
\]
Put $\sigma_j=z_j/|z_j|$, so $|\sigma_j|=1$ and $\sigma_j\in\{-1,1\}$ in the real case, and set
\[
  v_j=\overline{\sigma_j}\sqrt{\frac{a_j}{T}}w_j.
\]
Let $V$ have columns $v_j$ and put $X=V^*V$.  Then $X\succeq0$,
$\tr X=\sum_j\norm{v_j}^2=1$, and $X_{jj}=a_j/T$.  Moreover
\[
  Vz=\sum_jz_jv_j
  =\sum_j\frac{a_j}{\sqrt T}w_j=0.
\]
Thus $Xz=0$ and $\ran X\subseteq z^\perp=E_{\min}(G)$.  Since $q\ge2$ and $G$ has rank one,
$\lambda_{\min}(G)=0$.  Theorem~\ref{thm:energy-criterion} gives
$\nu(R)=T$ and the energy law is minimax.

We now prove the dominant branch directly.  Let $t=\pi_h$ and let
$w=z_{-h}/\sqrt b$ be the unit tail direction.  Multiplying $z$ by one global unit phase does not
change $G=zz^*$, so take $z_h=\sqrt a$ before forming the compression.  Compress $Q_\pi$ to
$\operatorname{span}\{e_h,w\}$.  Its first diagonal entry and off-diagonal magnitude are
\[
  \frac{a(1-t)}{t},\qquad \sqrt{ab}.
\]
For the second diagonal entry, Cauchy--Schwarz gives
\[
  \sum_{j\ne h}\frac{a_j^2}{b\pi_j}-b
  \ge\frac{b}{1-t}-b=\frac{bt}{1-t}.
\]
Consequently, by compression and Loewner monotonicity,
\begin{align*}
  \lambda_{\max}(Q_\pi)
  &\ge
  \lambda_{\max}
  \begin{pmatrix}
    a(1-t)/t&-\sqrt{ab}\\
    -\sqrt{ab}&bt/(1-t)
  \end{pmatrix}\\
  &=\frac{a(1-t)}{t}+\frac{bt}{1-t}
  \ge2\sqrt{ab}.
\end{align*}
The $2\times2$ matrix has determinant zero, which gives the equality on the second line.
The final inequality is sharp only when
$t=\sqrt a/(\sqrt a+\sqrt b)$.

Take the law in \eqref{eq:rank-one-law}.  Put
$d_h=a/\pi_h^*=a+\sqrt{ab}$ and
$d_T=a_j/\pi_j^*=b+\sqrt{ab}$ for every $j\ne h$.  On
$\operatorname{span}\{e_h,w\}$, the covariance is
\[
  \begin{pmatrix}
    \sqrt{ab}&-\sqrt{ab}\\
    -\sqrt{ab}&\sqrt{ab}
  \end{pmatrix},
\]
whose eigenvalues are $0$ and $2\sqrt{ab}$.  On the tail subspace orthogonal to $w$, the
eigenvalue is $b+\sqrt{ab}<2\sqrt{ab}$ because $a>b$.  Thus this law attains the lower bound.
Equality in the lower-bound proof forces both the unique value of $t$ and equality in the tail
Cauchy--Schwarz inequality, hence $\pi_j\propto a_j$ on the tail.  The dominant law is unique.

Finally, uniqueness in the balanced case follows from its positive-diagonal dual certificate.
For that certificate the reciprocal inner objective is strictly convex in $\pi$; every primal
optimum must minimize the same inner objective and therefore equals the energy law.
\end{proof}

\begin{remark}[The threshold]
At $a=T/2$, the polygon may close degenerately and the balanced formula remains valid.  Above
the threshold, one column cannot be cancelled by all the others.  The optimum then treats the
entire tail as one opposing block.  The phase transition is geometric, not numerical.
\end{remark}

\section{An unbounded separation from the Frobenius law}

\begin{corollary}[No universal multiplicative bound]
\label{cor:unbounded}
There is no finite constant $C$ such that the exact worst-query variance of squared-energy
sampling is at most $C\nu(R)$ for every residual $R$ with at least two nonzero columns, even under
the normalization $\norm{R}_F=1$.
\end{corollary}

\begin{proof}
Take two orthogonal columns with energies $a\ge b>0$ and $a+b=1$.  Under the squared-energy law
$(a,b)$, the two diagonal covariance entries are $b$ and $a$, so its exact worst-query variance is
$a$.  Theorem~\ref{thm:two-column} gives $\nu(R)=\sqrt{ab}$.  Their ratio is
\[
  \frac{a}{\sqrt{ab}}=\sqrt{\frac ab},
\]
which tends to infinity as $b\downarrow0$.
\end{proof}

The conclusion is stronger than a constant-factor sharpening of an existing estimate.  A
description chosen to minimize Frobenius residual energy need not minimize the exact fixed-law
access complexity.  If $\mathcal C_{d,m}(A)$ denotes a declared family of resident descriptions
of $A$ under description and metadata budgets $(d,m)$, the corresponding mathematical objective
is
\begin{equation}
  D_A^{\mathrm{mm}}(d,m)
  =\inf_{B\in\mathcal C_{d,m}(A)}\nu(A-B).
  \label{eq:minimax-distortion}
\end{equation}
Equation~\eqref{eq:minimax-distortion} is a proposed mathematical foundation, not a production
contract.  It says what would need to replace a Frobenius-only objective if exact worst-query
mean-square access were adopted.

\section{Physically fixed page components}

Columns are the clearest model, but a physical system may access aligned groups.  Let a declared
decomposition
\begin{equation}
  R=\sum_{g=1}^mR_g
  \label{eq:page-decomposition}
\end{equation}
be fixed before optimization, with every $R_g\ne0$.  On drawing $J\sim\pi$, return
$R_Jx/\pi_J$.

Put $A_g=R_g^*R_g$ and $B=R^*R$.  For $\pi\in\Delta_m^\circ$, define
\begin{equation}
  Q_\pi=\sum_{g=1}^m\frac{A_g}{\pi_g}-B
  \label{eq:page-covariance}
\end{equation}
and
\begin{equation}
  \nu(R_1,\ldots,R_m)=
  \min_{\pi\in\Delta_m^\circ}\lambda_{\max}(Q_\pi).
  \label{eq:page-primal}
\end{equation}

\begin{theorem}[Page-component dual]
\label{thm:page-dual}
The estimator has exact covariance $Q_\pi$, and the one-access minimax value is
\begin{equation}
  \nu(R_1,\ldots,R_m)
  =\max_{X\succeq0,\,\tr X=1}
  \left[
    \left(\sum_{g=1}^m\sqrt{\tr(XA_g)}\right)^2-\tr(XB)
  \right].
  \label{eq:page-dual}
\end{equation}
A primal law $\pi\in\Delta_m^\circ$ and value $t$ are optimal if and only if they admit a density
matrix $X$ with $\tr(XA_g)>0$ for every $g$ and
\begin{equation}
  Q_\pi\preceq tI,\quad X(tI-Q_\pi)=0,\quad
  \pi_g=\frac{\sqrt{\tr(XA_g)}}{\sum_h\sqrt{\tr(XA_h)}}.
  \label{eq:page-certificate}
\end{equation}
In particular, at least one positive-radicand certificate exists at every optimum; not every dual
optimizer must have this property.
\end{theorem}

\begin{proof}
Unbiasedness follows from \eqref{eq:page-decomposition}.  Expanding the second moment gives
\[
  \E\norm{R_Jx/\pi_J-Rx}^2
  =x^*\left(\sum_gA_g/\pi_g-B\right)x,
\]
which proves \eqref{eq:page-covariance} and its positive semidefiniteness.  The proof of
Theorem~\ref{thm:dual} now applies word for word with
$a_jX_{jj}$ replaced by $\tr(XA_g)$.  Coercivity remains valid because
\[
  \lambda_{\max}(Q_\pi)
  \ge \frac{\lambda_{\max}(A_g)}{\pi_g}-\norm{B}_2
\]
for each nonzero $A_g$.  The sufficiency proof of Theorem~\ref{thm:certificate} gives the reverse
bound from any witness satisfying \eqref{eq:page-certificate}.

For necessity, a saddle-point witness obeys
$\tr(XA_g)/\pi_g^2=\gamma$ independently of $g$.  If the optimal value $t$ is positive, then
$\gamma>0$: otherwise $R_gX^{1/2}=0$ for every $g$, hence $RX^{1/2}=0$ and
$t=\tr(XQ_\pi)=0$, a contradiction.  Thus every radicand is positive and normalization gives
\eqref{eq:page-certificate}.  If $t=0$, the covariance identity is a sum of nonnegative squared
errors and forces $R_g=\pi_gR$ for every $g$.  Choose a unit vector $u$ with $Ru\ne0$ and set
$X=uu^*$.  Then $Q_\pi=0$ and
$\tr(XA_g)=\pi_g^2\norm{Ru}_2^2>0$, so the same certificate equation holds.
\end{proof}

\begin{theorem}[Zero-variance page decompositions]
\label{thm:page-zero}
The value $\nu(R_1,\ldots,R_m)$ is zero if and only if there is a probability vector
$\pi\in\Delta_m^\circ$ such that
\begin{equation}
  R_g=\pi_gR\qquad(1\le g\le m).
  \label{eq:proportional-pages}
\end{equation}
\end{theorem}

\begin{proof}
For every $x$,
\begin{align*}
  x^*Q_\pi x
  &=\sum_g\pi_g
  \norm{R_gx/\pi_g-Rx}_2^2.
\end{align*}
If \eqref{eq:proportional-pages} holds, every summand vanishes.  Conversely, if the minimax value
is zero, attainment gives a law with $Q_\pi=0$.  The displayed sum then vanishes for every $x$.
Every term is nonnegative and every $\pi_g$ is positive, so
$R_gx/\pi_g=Rx$ for every $g$ and every $x$, which is \eqref{eq:proportional-pages}.
\end{proof}

The decomposition in \eqref{eq:page-decomposition} must be physical data, not a free algebraic
choice.  Otherwise one could split $R$ into proportional replicas and force the invariant to
zero without changing the actual access mechanism.

\section{Boundaries of the theorem}

The invariant is homogeneous of degree two.  It is unchanged by output isometries, column
permutations, and coordinatewise input phases.  It is not invariant under a general right-unitary
change of basis, nor should it be: the addressable coordinates are part of the information model.

In the column model, $\nu(R)=0$ if and only if $R$ has exactly one nonzero column.  Indeed,
$Q_\pi=0$ forces the off-diagonal entries of $G$ to vanish and then forces $\pi_j=1$ for every
nonzero column, possible only for one column.  The page model is less rigid because distinct page
operators can be proportional to the same $R$.

Allowing a different independent law at each of $s$ times does not improve the ordinary mean.
If the laws are $\pi^{(1)},\ldots,\pi^{(s)}$ and
$\bar\pi=s^{-1}\sum_t\pi^{(t)}$, convexity of $u\mapsto1/u$ gives
\[
  \frac1s\sum_{t=1}^sQ_{\pi^{(t)}}\succeq Q_{\bar\pi}.
\]
Therefore the covariance of their arithmetic mean has largest eigenvalue at least $\nu/s$.
The common optimal law already attains $\nu/s$.

Sampling without replacement can do better.  For $R=I_2$, two deterministic distinct column
accesses reconstruct $Rx$ exactly, while two independent categorical samples have minimax
worst-query mean-square error $\nu(I_2)/2=1/2$.  No theorem above claims otherwise.

\section{Computational falsification}

A standard-library program accompanies the manuscript.  It is a counterexample search, not a
proof.  Its fixed-seed checks compare:
\begin{enumerate}[label=\textup{(\alph*)},leftmargin=2.1em]
  \item direct enumeration of sampling outcomes with the covariance quadratic form;
  \item the two-column closed form with a primal simplex grid and a signed rank-one
  density-witness boundary grid;
  \item the asserted and reported orthogonal equalization residual and a three-coordinate simplex
  search; and
  \item both branches of the rank-one formula with a three-coordinate simplex grid.
\end{enumerate}
The search covers 400 covariance identities, 160 two-column residuals, 30 orthogonal triples,
and 80 rank-one triples, for 670 fixed-seed cases, plus the embedded grid comparisons.  The primal
grid denominators are 600, 120, and 160; the dual-boundary denominator is 4000.  These finite
grids give no off-grid optimality guarantee.  The covariance, primal-grid, and closed-form
evaluations share a Gram helper and a floating Jacobi eigensolver; the dual-boundary and
equalization checks are separate scalar computations.  Agreement is diagnostic only.  The
program also illustrates Corollary~\ref{cor:unbounded} at seven finite ratios; the limiting claim is
proved above.  Numerical agreement cannot replace any proof above; disagreement would falsify the
corresponding formula or the program.

One boundary error was found during this process.  An early statement assigned zero accesses
whenever $\nu(R)=0$.  A one-column nonzero residual disproved it: its estimator has zero variance
after one exact access, but zero accesses do not evaluate the residual.  The corrected distinction
appears in Theorem~\ref{thm:sample-complexity}.

\section{Consequences for Cassette}

The mathematical consequence is exact and conditional.  If Cassette retains the declared
query-independent, with-replacement estimator, then $\nu$ is the necessary and sufficient
worst-unit-query mean-square constant.  The present squared-energy law is exact precisely under
Theorem~\ref{thm:energy-criterion}; outside that geometry, it can be strictly and even
unboundedly worse.

A later adoption would require more than replacing one formula.  The resident-description
objective would become \eqref{eq:minimax-distortion}; a plan would carry the physical decomposition;
and a certificate would carry the primal law, value, and density witness.  Those are mathematical
implications, not implementation instructions.

Nothing in this paper is integrated into Cassette production.  The production mathematics,
schemas, plans, tests, and stage status remain unchanged.  The result is presented for separate
review before any adoption decision.

\section{Conclusion}

The fixed-law categorical residual problem has an exact spectral geometry.  Its covariance is a
reciprocal diagonal perturbation of a Gram matrix; its minimax value is dual to a density-matrix
balancing problem; and its optimal law can differ without bound from squared-energy sampling.
The dual resolves repeated eigenvalues, the closed forms expose the controlling geometry, and the
page extension ties the invariant to a physical decomposition.  The narrow result is a candidate
foundation for exact worst-query residual-access certificates.  Whether Cassette adopts that
foundation is a separate engineering and governance decision.

\begin{thebibliography}{99}

\bibitem{HansenHurwitz1943}
M. H. Hansen and W. N. Hurwitz,
``On the theory of sampling from finite populations,''
\emph{Annals of Mathematical Statistics}, vol.~14, no.~4, pp.~333--362, 1943.
\href{https://doi.org/10.1214/aoms/1177731356}{doi:10.1214/aoms/1177731356}.

\bibitem{DrineasKannanMahoney2006}
P. Drineas, R. Kannan, and M. W. Mahoney,
``Fast Monte Carlo algorithms for matrices I: Approximating matrix multiplication,''
\emph{SIAM Journal on Computing}, vol.~36, no.~1, pp.~132--157, 2006.
\href{https://doi.org/10.1137/S0097539704442684}{doi:10.1137/S0097539704442684}.

\bibitem{ErikssonBiqueEtAl2011}
S. Eriksson-Bique, M. Solbrig, M. Stefanelli, S. Warkentin, R. Abbey, and I. C. F. Ipsen,
``Importance sampling for a Monte Carlo matrix multiplication algorithm, with application to
information retrieval,''
\emph{SIAM Journal on Scientific Computing}, vol.~33, no.~4, pp.~1689--1706, 2011.
\href{https://doi.org/10.1137/10080659X}{doi:10.1137/10080659X}.

\bibitem{HolodnakIpsen2015}
J. T. Holodnak and I. C. F. Ipsen,
``Randomized approximation of the Gram matrix: Exact computation and probabilistic bounds,''
\emph{SIAM Journal on Matrix Analysis and Applications}, vol.~36, no.~1, pp.~110--137, 2015.
\href{https://doi.org/10.1137/130940116}{doi:10.1137/130940116}.

\bibitem{AchlioptasMcSherry2007}
D. Achlioptas and F. McSherry,
``Fast computation of low-rank matrix approximations,''
\emph{Journal of the ACM}, vol.~54, no.~2, article 9, 2007.
\href{https://doi.org/10.1145/1219092.1219097}{doi:10.1145/1219092.1219097}.

\bibitem{ImbergEtAl2023}
H. Imberg, M. Axelson-Fisk, and J. Jonasson,
``Optimal subsampling designs,'' arXiv:2304.03019, 2023.
\href{https://arxiv.org/abs/2304.03019}{arXiv:2304.03019}.

\bibitem{Kiefer1974}
J. Kiefer,
``General equivalence theory for optimum designs (approximate theory),''
\emph{Annals of Statistics}, vol.~2, no.~5, pp.~849--879, 1974.
\href{https://doi.org/10.1214/aos/1176342810}{doi:10.1214/aos/1176342810}.

\bibitem{HarmanRosa2018}
R. Harman and S. Rosa,
``Removal of the points that do not support an E-optimal experimental design,''
arXiv:1808.00731v1, 2018.
\href{https://arxiv.org/abs/1808.00731v1}{arXiv:1808.00731v1}.

\bibitem{Sion1958}
M. Sion,
``On general minimax theorems,''
\emph{Pacific Journal of Mathematics}, vol.~8, no.~1, pp.~171--176, 1958.
\href{https://doi.org/10.2140/pjm.1958.8.171}{doi:10.2140/pjm.1958.8.171}.

\bibitem{Horn1954}
A. Horn,
``Doubly stochastic matrices and the diagonal of a rotation matrix,''
\emph{American Journal of Mathematics}, vol.~76, no.~3, pp.~620--630, 1954.
\href{https://doi.org/10.2307/2372705}{doi:10.2307/2372705}.

\bibitem{CasazzaEtAl2006}
P. G. Casazza, M. Fickus, J. Kova\v{c}evi\'c, M. T. Leon, and J. C. Tremain,
``A physical interpretation of tight frames,'' in
\emph{Harmonic Analysis and Applications}, pp.~51--76, Birkh\"auser, 2006.
\href{https://doi.org/10.1007/0-8176-4504-7_4}{doi:10.1007/0-8176-4504-7\_4}.

\bibitem{Novak1992}
E. Novak,
``Optimal linear randomized methods for linear operators in Hilbert spaces,''
\emph{Journal of Complexity}, vol.~8, no.~1, pp.~22--36, 1992.
\href{https://doi.org/10.1016/0885-064X(92)90032-7}{doi:10.1016/0885-064X(92)90032-7}.

\bibitem{RavichandranEtAl2024}
A. Ravichandran, N. E. Pashley, B. Libgober, and T. Dasgupta,
``Optimal allocation of sample size for randomization-based inference from $2^K$ factorial
designs,''
\emph{Journal of Causal Inference}, vol.~12, no.~1, article 20230046, 2024.
\href{https://doi.org/10.1515/jci-2023-0046}{doi:10.1515/jci-2023-0046}.

\end{thebibliography}

\end{document}
